\section{Enumeration by reversible finite constructions}
\label{sec:enumeration}
The historical example of six tastes, the twelvefold way, and the five
binary trees on three vertices all occur in the seed article. Its links to
\emph{Binomial coefficient}, \emph{Composition (combinatorics)},
\emph{Fibonacci sequence}, \emph{Partition of a set}, \emph{Catalan number},
and \emph{Schr\"oder--Hipparchus number} lead to the following calculations.
Write $[n]=\{0,\ldots,n-1\}$, and give a word of fixed length its positional
code. Throughout, a bijection means two specified finite algorithms whose
composites are the identity on the indicated valid codes. Counting equality
then follows by summing the graph of either algorithm in the two orders.

\subsection{Subsets, lists, and compositions}
Let $b(n,k)$ count length-$n$ bit strings of weight $k$. Thus $b(n,k)=0$
for $k>n$. We also write $\binom nk=b(n,k)$ and put $0^0=1$.
\begin{theorem}[Binomial identities]\label{result04}
For $k\le n$, $b(n,k)=b(n,n-k)$ and
$b(n,k)k!(n-k)!=n!$. Moreover
\begin{align*}
 b(n+1,k+1)&=b(n,k)+b(n,k+1),\\
 \sum_{j=0}^k b(m,j)b(n,k-j)&=b(m+n,k),\\
 (a+b)^n&=\sum_{k=0}^n b(n,k)a^k b^{n-k},\\
 \sum_{j=k}^n b(j,k)&=b(n+1,k+1) \quad(k\le n),\\
 k b(n,k)&=n b(n-1,k-1)\quad(1\le k\le n).
\end{align*}
In particular the row sum is $2^n$, the sum of the squares in row $n$ is
$b(2n,n)$, and there are $63$ nonempty choices from six tastes.
\end{theorem}
\begin{proof}
Complement each bit for symmetry. Sort the first $k$ entries of a
permutation into a set, retaining the two orderings on this set and its
complement. Putting the two lists back to back reverses the operation.
The number of permutations is $n!$, since the position of the largest
entry has $n$ possibilities after the other entries are permuted. Separating
the last bit gives Pascal's rule. Cutting after $m$ bits gives the second
identity, with $j$ the weight in the first piece. A word on the disjoint
union of alphabets of sizes $a,b$ is recovered from the positions occupied by
the first alphabet and the two extracted words. A nonempty $(k+1)$-set
in $[n+1]$ has a largest element $j$ and $k$ other elements in $[j]$.
Finally a $k$-set with a marked member is determined by that member and a
$(k-1)$-set in its ordered complement. The square sum follows by splitting
a length-$2n$ string of weight $n$ at its middle and complementing the
second piece. All relabelings here and below use increasing order.
\end{proof}

The sequence of binomial coefficients is unimodal. Indeed,
$(k+1)b(n,k+1)=(n-k)b(n,k)$ by the same marked-subset identity (or by
forming a $k$-set and adding a new element). Thus $b(n,k+1)\ge b(n,k)$
when $k+1\le n-k$, and the reversed inequality holds on the other half.
For ordered words with prescribed multiplicities $a_1,\ldots,a_r$ summing
to $n$, let $M$ be their count. Distinguishing occurrences of each letter
by a chosen linear order gives a permutation of $n$ labeled items, so
$M\prod_i a_i!=n!$. This also proves the multinomial expansion by counting
words on a disjoint union of $r$ colored alphabets.

\begin{proposition}[Compositions and Fibonacci identities]\label{result05}
The positive compositions of $n>0$ into $r>0$ parts number $b(n-1,r-1)$;
there are $2^{n-1}$ in all. For $r>0$, nonnegative $r$-tuples summing to
$n$ number $b(n+r-1,r-1)$. With $F_0=0,F_1=1,F_{t+2}=F_{t+1}+F_t$,
\[
 F_{n+1}=\sum_{j=0}^{\lfloor n/2\rfloor} b(n-j,j),\qquad
 F_{m+n+1}=F_{m+1}F_{n+1}+F_m F_n .
\]
Also $\sum_{i=0}^n F_i=F_{n+2}-1$,
$\sum_{i=0}^n F_i^2=F_nF_{n+1}$ and
$F_{n+1}F_{n-1}-F_n^2=(-1)^n$ for $n\ge1$.
\end{proposition}
\begin{proof}
Choose $r-1$ cuts among the $n-1$ gaps in a row of $n$ marks, or allow
any gap set if the number of parts is free. For a nonnegative tuple write
$n$ marks with $r-1$ bars, permitting adjacent or terminal bars. Reading
the run lengths inverts the map. For the Fibonacci identities, tile a
length-$n$ strip by monomers and dominoes. The empty strip and the unit
strip each have one tiling. Removing the last tile proves by induction that
the count is $F_{n+1}$. With $j$ dominoes there are $n-j$ tiles, and one
chooses the $j$ domino positions. At a given seam after cell $m$ in a
strip of length $m+n$, either the two strips are separately tiled or one
domino crosses the seam. The crossing case leaves lengths $m-1,n-1$,
and has count zero at an endpoint, as prescribed by $F_0=0$. The first
sum telescopes using $F_i=F_{i+2}-F_{i+1}$, and the sum of squares
telescopes using $F_iF_{i+1}-F_{i-1}F_i=F_i^2$ for $i\ge1$.
Finally $D_n=F_{n+1}F_{n-1}-F_n^2$ has $D_1=-1$ and
$D_{n+1}=F_{n+1}(F_n-F_{n+1})+F_n^2=-D_n$. Signed identities can
equivalently be stated by moving negative terms to the opposite side.
\end{proof}

\subsection{All twelve cases for distributing objects}
We distinguish positions (the $n$ objects) and labels (the $r$ boxes).
Put $(r)_n=r(r-1)\cdots(r-n+1)$ for $n\le r$, and zero otherwise.
Put $W(n,r)=b(n+r-1,r-1)$ and $V(n,r)=[n\ge r]b(n-1,r-1)$
when $r>0$ (the second is zero for $n=0$). For $r=0$ put
$W(n,0)=V(n,0)=[n=0]$. Here $[P]$ is $1$ or $0$ according as the
decidable test $P$ succeeds or fails. Write $S(n,j)$ for the number of
partitions of $[n]$ into $j$ nonempty blocks. Write $p_j(n)$ for the
number of nonincreasing lists of $j$ positive integers of sum $n$, and
$p_{\le r}(n)=\sum_{j=0}^r p_j(n)$. The empty partition contributes $1$.

\begin{theorem}[Twelvefold way, including empty sets]\label{result06}
The following are the numbers of classes of functions $[n]\to[r]$.
A label declared immaterial may be permuted arbitrarily.
\[
\begin{array}{ll|ccc}
\text{objects}&\text{boxes}&\text{all}&\text{injective}&\text{surjective}\\ \hline
\text{labeled}&\text{labeled}&r^n&(r)_n&r!S(n,r)\\
\text{unlabeled}&\text{labeled}&W(n,r)&b(r,n)&V(n,r)\\
\text{labeled}&\text{unlabeled}&\sum_{j=0}^r S(n,j)&[n\le r]&S(n,r)\\
\text{unlabeled}&\text{unlabeled}&p_{\le r}(n)&[n\le r]&p_r(n)
\end{array}
\]
\end{theorem}
\begin{proof}
A labeled function is a word. For an injection the next position can take
each unused label once. For a surjection take its nonempty fibers, order
these blocks by their least elements, and record the permutation assigning
the $r$ labels to them. These choices reconstruct the word.
Ignoring object labels leaves exactly the vector of occupancy numbers:
match the increasing lists of objects in corresponding fibers to obtain
the relabeling. The vector is a weak composition, a bit vector of weight
$n$, or a positive composition, respectively. Ignoring just box labels
leaves exactly the equality partition on positions. A canonical word uses
the least-element order of its nonempty fibers for consecutive labels;
the remaining labels can be matched in increasing order. The injection
then has only the singleton partition when $n\le r$.
Ignoring both kinds of label leaves the sorted positive occupancies.
The matching algorithm first matches equal-sized boxes in increasing order
within each size (including size zero) and then matches their ordered
contents. Injectivity leaves only a list of $n$ ones if $n\le r$.
The map with empty domain is unique for every $r$ and is onto precisely
when $r=0$; a positive domain admits no map to an empty set.
\end{proof}

\begin{corollary}[Stirling and Bell identities]\label{result07}
Let $B_n=\sum_{j=0}^n S(n,j)$. Then $S(0,0)=1$, $S(n+1,0)=0$, and
\[
\begin{split}
 S(n+1,k+1)&=S(n,k)+(k+1)S(n,k+1),\\
 B_{n+1}&=\sum_{j=0}^n b(n,j)B_j,\qquad
 r^n=\sum_{j=0}^n S(n,j)(r)_j .
\end{split}
\]
\end{corollary}
\begin{proof}
Delete the new element $n$. If its block disappears, there remain $k$
blocks; otherwise the old partition has $k+1$ blocks and one records
which received $n$. To count all partitions choose the $j$ elements outside
the block of $n$ and partition them; all the remaining elements join $n$.
For the last identity partition the positions by their values and assign
distinct labels to the blocks in canonical order.
\end{proof}

For completeness the companion numbers $c(n,k)$ count permutations of
$[n]$ with exactly $k$ cycles. Cycles can be computed by following each
unvisited least entry until its first return (at most $n$ steps). A repeated
entry must be the starting point, since the predecessor is unique. Removing
$n$ gives either a singleton cycle or a choice of one of the $n$ old
elements after which to reinsert $n$. Thus
$c(n+1,k+1)=c(n,k)+n c(n,k+1)$ with the same boundary row as $S$.
Since multiplication of a polynomial by $x+n$ has this coefficient
recursion, $\sum_k c(n,k)x^k=x(x+1)\cdots(x+n-1)$.
Replacing $x$ by $-x$ and taking the overall sign gives
$(x)_n=\sum_k(-1)^{n-k}c(n,k)x^k$ as a polynomial over signed numbers.
Conversely $x(x)_j=(x)_{j+1}+j(x)_j$ and the recursion for $S$ give the
polynomial identity $x^n=\sum_jS(n,j)(x)_j$ by induction on $n$.
Substitution and coefficient comparison give the triangular inverse identity
$\sum_{j=k}^n S(n,j)(-1)^{j-k}c(j,k)=[n=k]$.
All comparisons take place in arrays of degree at most $n$.

\subsection{Inclusion--exclusion and permutations}
\begin{theorem}[A cancellation involution]\label{result08}
For a finite universe $U$ and a coded family $A_0,\ldots,A_{m-1}$ of
subsets, put $I_T=\#\{x\in U:\forall i\in T,\ x\in A_i\}$. Then
\[
 \#\{x\in U:\forall i<m,\ x\notin A_i\}
 =\sum_{T\subseteq[m]}(-1)^{|T|}I_T .
\]
In particular, for the number $D_n$ of fixed-point-free permutations and
the number of surjections one obtains
\[
 D_n=\sum_{j=0}^n(-1)^j b(n,j)(n-j)!,
 \qquad r!S(n,r)=\sum_{j=0}^r(-1)^j b(r,j)(r-j)^n .
\]
\end{theorem}
\begin{proof}
For fixed $x$ form $H=\{i:x\in A_i\}$. In the signed sum its occurrences
are indexed exactly by $T\subseteq H$. If $H$ is empty, $T$ is empty and
gives one positive occurrence. Otherwise toggling its least member pairs
each positive occurrence with one negative occurrence. In a permutation
fixing a specified $j$-set, the complement may be freely permuted. In a
word avoiding $j$ specified labels, each position has $r-j$ choices.
Using as bad events fixed points and missing labels respectively proves the
formulas. In particular both expressions are identities of natural
numbers once the odd-indexed terms are transposed.
\end{proof}
Choosing the fixed set also gives $n!=\sum_{j=0}^n b(n,j)D_{n-j}$.
For $n\ge2$ the first value $j$ of a derangement of $[n]$ has $n-1$
choices. If $0$ and $j$ form a two-cycle, delete them. Otherwise delete
$0$ from its cycle, replacing the unique arrow into $0$ by an arrow into
$j$; the remaining permutation is a derangement. In the latter case the
predecessor of $j$ determines where to reinsert $0$. Hence
$D_n=(n-1)(D_{n-1}+D_{n-2})$, with $D_0=1,D_1=0$.

\subsection{Integer partitions as finite coefficient identities}
For each degree bound $N$ a generating function below means its first
$N+1$ coefficients. The coefficient of $x^t$ in a product is the sum
over $i+j=t$ of the two coefficients multiplied. This counts pairs of
objects of total size $t$ and gives associativity by summing over triples.
If $A$ counts positive-size objects, $1+A+\cdots+A^N$ counts sequences of
such objects through degree $N$, since a sequence of size $t$ has length
at most $t$. We abbreviate this finite geometric sum by $(1-A)^{-1}$.
This interpretation permits the usual coefficientwise series notation over
either natural or signed numbers.

\begin{theorem}[Conjugation, Euler's product, and the Durfee square]\label{result09}
Write $p(n)=p_{\le n}(n)$, $P_d(x)=\prod_{j=1}^d(1-x^j)^{-1}$.
Then $[x^n]P_d(x)=p_{\le d}(n)$ and
\[
 p(n)=[x^n]P_n(x)
     =[x^n]\sum_{d=0}^{\lfloor\sqrt n\rfloor}x^{d^2}P_d(x)^2 .
\]
For $n\ge d$, $p_d(n)=p_{\le d}(n-d)$, also with $d=0$.
For every integer $h\ge2$, the number of partitions of $n$ having no
part divisible by $h$ equals the number having every multiplicity $<h$.
In particular partitions into odd parts and into distinct parts are
equinumerous.
\end{theorem}
\begin{proof}
A partition drawn in left-justified rows of lengths $\lambda_i$ is encoded
by a decreasing list padded with zeros. Its transposed row $j$ has length
$\#\{i:\lambda_i\ge j\}$ for $j\ge1$. Cellwise, $(i,j)$ belongs to
the transpose exactly when $(j,i)$ belongs to the original shape. Thus
transposition interchanges maximum part and number of nonzero parts, and
preserves size. The product $P_d$ selects each multiplicity of parts of
sizes $1$ through $d$, so counts partitions with largest part at most $d$.
Subtracting one from each of $d$ positive rows gives a partition of $n-d$
with at most $d$ rows; pad to length $d$ before adding back.
For a size-$n$ partition take the greatest $d$ with $\lambda_d\ge d$
(using $0$ as an always admissible value). Remove its northwest $d$ by $d$
square. The cells to the right form a partition with at most $d$ rows;
those below have largest part at most $d$. Conversely these two shapes
adjoin uniquely to the square, and their bounds ensure no larger square
fits. This proves the coefficient formula, including the empty diagram.

Finally for each $j$ not divisible by $h$, write the multiplicity $m_j$
of $j$ in base $h$, say $m_j=\sum_t a_t h^t$. Replace the $m_j$ copies by
$a_t$ copies of $jh^t$ for each $t$. Every positive integer $v$ has a
unique expression $jh^t$ with $h\nmid j$: repeatedly divide by $h$ while
divisible, at most $v$ times. Thus output multiplicities are $<h$, weight
is preserved, and splitting every part $jh^t$ into $h^t$ copies of $j$
gives the inverse. All loops on positive-weight parts and nonzero digits
are bounded by $n$.
\end{proof}

\subsection{Catalan objects and Schr\"oder dissections}
A Dyck word of semilength $n$ has $n$ opening and $n$ closing steps, with
every prefix containing at least as many opens as closes. Let $C_n$ count
these words. A step changes the signed height by one.
\begin{theorem}[Reflection and first return]\label{result10}
We have $C_0=1$ and
\[
 C_n=b(2n,n)-b(2n,n+1),\qquad
 (n+1)C_n=b(2n,n),\qquad
 C_{n+1}=\sum_{j=0}^n C_j C_{n-j}.
\]
Thus $C(x)=1+xC(x)^2$ coefficientwise.
\end{theorem}
\begin{proof}
A balanced word that fails its prefix test first hits $-1$ at a bounded
position. Swap the two letters up to and including this position. The
new word has $n+1$ opens, total height two, and its first hit of $+1$ is
the changed position. Any length-$2n$ word with $n+1$ opens reaches $+1$
because its height starts at zero, ends at two and advances by unit steps.
Reflect there to invert the map. For $n=0$ both exceptional classes are
empty. The marked-subset relation gives
$(n+1)b(2n,n+1)=n b(2n,n)$, including $n=0$. Multiply the first identity
by $n+1$ and cancel to obtain the second. A nonempty Dyck word splits
uniquely into an open, a Dyck word, a close, and a second Dyck word:
the displayed close is the first positive-time return to height zero.
Concatenating reverses the split and proves the convolution.
\end{proof}

\begin{proposition}[Concrete Catalan families]\label{result11}
$C_n$ counts full ordered binary trees with $n$ internal vertices,
binary tree shapes with $n$ vertices and distinguished left and right
slots, full parenthesizations of $n+1$ ordered letters, plane rooted trees
with $n$ edges, noncrossing pairings on $[2n]$, noncrossing partitions of
$[n]$, two-row standard tableaux of shape $(n,n)$, and, for $n\ge1$,
triangulations of a labeled convex $(n+2)$-gon in its fixed cyclic order.
In particular $C_n\le B_n$, and $C_3=5$.
\end{proposition}
\begin{proof}
Make an empty Dyck word a leaf, and make the two words in the first-return
decomposition the ordered children of an internal vertex. Recursive
expansion stops after at most $n$ internal vertices. Deleting external
leaves and retaining empty child slots, or writing an operation at each
internal vertex and assigning letters to the leaves in order, gives the
next two forms. There are $n+1$ leaves by induction on this construction.
For a plane tree record an open on descending each edge during the
left-to-right traversal, and a close on ascending. Conversely start at a
root, add a new last child and descend on each open, and ascend on each
close. Height is depth, so every step is valid and the walk ends at the
root. For a pairing, pair each close with the most recent unmatched open.
These nested or disjoint pairs cannot cross. Conversely writing opens at
the smaller endpoints of a noncrossing pairing yields this same matching:
skipping a more recent unmatched open would force its eventual pair to
cross the chosen pair. For the tableau, place $i=1,\ldots,2n$ in the next
free column of the top or bottom row according to the $i$th letter.
The prefix condition is exactly the condition that the top entry of each
column is filled before its bottom entry.

For a noncrossing partition of $[n]$, let $r_i$ be the size of the block
whose least member is $i$, or zero if there is none. At each $i$ print
$r_i$ opens followed by one close. At every stage the blocks started so
far contain all elements seen so far; the prefix condition follows.
Conversely a Dyck word uniquely splits at its $n$ closes into these runs.
When $r_i>0$, push a new block with $r_i$ vacant slots. Assign $i$ to
the top block, decrease its vacancies and pop it when full. Stack height
in \emph{slots} is the word height, so the assignment exists and all
blocks finish. An older block pauses until every block above it finishes,
which excludes an interlacing $a<b<c<d$ belonging alternately to two
different blocks. Conversely, if an old block received $i$ while a younger
one was still active, their least elements, $i$, and a future member of
the younger block would interlace. This proves both inverses and the
claimed inequality by forgetting the noncrossing test.

In a triangulation the triangle incident to the distinguished base edge
has a unique third vertex. Its other two sides bound two smaller
triangulated intervals of boundary vertices, allowing a bare edge for an
empty interval. Conversely, glue such a pair to the base triangle.
The positions of the third vertex correspond to $j=0,\ldots,n-1$, and
the two child objects have $j$ and $n-1-j$ triangles. This is exactly
the binary tree construction, with a bare edge representing a leaf.
Finally the convolution gives $C_0,C_1,C_2,C_3=1,1,2,5$.
\end{proof}

\begin{proposition}[Schr\"oder--Hipparchus numbers]\label{result12}
Let $H_n$ count ordered rooted trees with $n\ge1$ leaves and at least two
children at every internal vertex, with $H_0=0$. With $H=\sum H_nx^n$,
\[
 H=x+\frac{H^2}{1-H},\qquad
 H_1=1,\qquad H_n=2\sum_{i=1}^{n-1}H_iH_{n-i}-H_{n-1}\quad(n\ge2).
\]
These count unrestricted parenthesizations of $n$ ordered letters into
groups of arity at least two, with the outer pair omitted, and dissections
of a convex $(n+1)$-gon by any noncrossing subset of diagonals ($n\ge2$).
In particular $H_{10}=103049$.
\end{proposition}
\begin{proof}
A leaf contributes $x$. Otherwise the root carries a list of $r\ge2$
positive-size trees, so the contribution of all lists is
$H^2(1+H+\cdots)$, finite at each degree. Multiply by $1-H$ and compare
coefficients to get the recurrence. With the leaf base, it gives
$1,1,3,11,45,197,903,4279,20793,103049$ for leaf counts $1$ through $10$.
Recording child lists by grouping parenthesizes the letters, and labeled
letters are read left to right to recover the ordered tree.
In a polygon dissection start at the cell incident to the distinguished
base. Its other $r\ge2$ sides, in boundary order, are children. An outer
side is a leaf and an inner diagonal is the base of the dissection beyond
it. The inverse allocates consecutive boundary intervals of lengths equal
to the leaf counts of the children and uses their endpoints as vertices
of the base cell. This produces exactly the original diagonals, since
each inner side separates its interval from the rest. Restriction to
triangles recovers the binary case.
\end{proof}

\subsection{Narayana numbers and dissections with prescribed size}
The partition and tree bijections preserve enough information to refine
the preceding counts. A peak is an open immediately followed by a close.
For $m\ge1$ and $1\le k\le m$, let $N(m,k)$ count Dyck words of
semilength $m$ with $k$ peaks.
\begin{theorem}[Narayana refinement]\label{result13}
The number of noncrossing partitions of $[m]$ with $k$ blocks is $N(m,k)$,
where
\[
 mN(m,k)=b(m,k)b(m,k-1).
\]
Furthermore $H_{m+1}=\sum_{k=1}^{m}N(m,k)2^{k-1}$. For a convex
$(m+2)$-gon the number $L(m,d)$ of dissections with exactly $d$ diagonals
satisfies
\[
 (d+1)L(m,d)=b(m-1,d)b(m+d+1,d)\qquad(0\le d<m).
\]
\end{theorem}
\begin{proof}
A block in the partition encoding starts a nonempty run of opens ending
at a close. Thus blocks are counted by peaks. With $k$ peaks a Dyck word
has alternating positive up-run and down-run lengths summing to $m$ in
each direction. Choose the $k-1$ internal partial sums $A$ of the up-run
lengths and $B$ of the down-run lengths from $\{1,\ldots,m-1\}$.
In increasing order they satisfy $B_i\le A_i$, since after the $i$th
down-run total closes do not exceed total opens. Equivalently every
prefix of $1,\ldots,m-1$ has at least as many members of $B$ as of $A$.
Conversely taking successive differences, including the last gap up to
$m$, recovers the word. Put $q=m-1$ and $s=k-1$.
For a bad pair of $s$-sets the first prefix with one excess member of
$A$ over $B$ exists. Swap the sets on this prefix, to obtain sizes
$s-1$ and $s+1$, respectively. The inverse uses the first prefix having
more members of the new second set. It exists because the final
difference is two and changes by at most one per position. Consequently
\[
 N(m,k)= b(q,s)^2-b(q,s-1)b(q,s+1),
\]
where an out-of-range index gives zero. At either endpoint $s=0,q$
the claimed formula follows directly. Otherwise divide the last expression
by $b(q,s)^2$ in positive rational arithmetic. The ratio is
\[
 1-\frac{s(q-s)}{(q-s+1)(s+1)}
 =\frac{m}{(m-k+1)k}.
\]
Since $b(m,k)=m b(q,s)/k$ and
$b(m,k-1)=m b(q,s)/(m-k+1)$, this gives the natural-number identity
after clearing the positive denominators.

Replace each ordered node of arity $u\ge2$ by its binary right comb:
the first child is the left child, and subsequent children are attached
along successive right edges. Mark the $u-2$ new internal right edges.
Conversely in any full binary tree mark any subset of edges whose right
child is internal. The marked-edge components are rightward chains, so
contracting each chain to one vertex, retaining the left-to-right order
of its hanging children, inverts the expansion. A binary tree with $m$
internal vertices has $m-1$ edges between internal vertices. Under the
first-return encoding a peak means exactly that its vertex's left child
is external. Hence $m-k$ internal edges go left and $k-1$ go right. For
each of the $N(m,k)$ shapes there are $2^{k-1}$ choices.

Contracting $t$ edges leaves $m-t$ internal vertices, hence $m-1-t$
diagonals in the dissection: one diagonal per nonroot internal vertex.
Put $d=m-1-t$. Choosing exactly $t$ edges gives
\[
 \begin{split}
 m L(m,d)
  &=\sum_{k=1}^m b(m,k)b(m,k-1)b(k-1,t)\\
  &=b(m,t)\sum_{k=1}^m b(m,k)b(m-t,m+1-k)
    =b(m,t)b(2m-t,m+1).
 \end{split}
\]
Here a nested $t$-set in a $(k-1)$-set may be chosen first, and the
last equality splits an $(m+1)$-set between two rows. Substituting $t$,
using symmetry, and applying $(d+1)b(m,d+1)=m b(m-1,d)$ proves the
formula by cancellation. Each chain contraction and its inverse use at
most $2m+1$ vertices. The proof uses quotients only with explicit
positive denominators and has the displayed division-free output.
\end{proof}

\subsection{Arithmetic implementation and Lean proof obligations}
The following specifications give finite loop bounds for the less direct
bijections. They also allow the various identities to share a single Lean
counting lemma. These are targets for formalization.
\begin{itemize}
\item Define \texttt{wordOK} by extracting $n$ digits in base
$\max(2,r)$, testing all digits $<r$ and requiring the code to be less
than the $n$th power of the base. Bits give the subset case.
Define \texttt{kernelWord} by replacing each digit by the index of its
first occurrence among the distinct digits seen so far. A restricted
growth word starts at zero when nonempty, and each next label is at most
one plus the preceding maximum. Prove by prefix induction that the
normalized word has the same equality relation on positions and is fixed
by normalization. Its number of distinct letters is the number of blocks.
Normalized words for partitions of $[n]$ fit in $\max(2,n)^n$. The original word over $[r]$ fits in $\max(2,r)^n$. Empty domain and empty alphabet pass the general code test
exactly in the cases used above.
\item Define \texttt{occupancy} by counting matching digits, and
\texttt{representative} by concatenating the indicated repetitions.
Use increasing rank within each fiber to construct the permutations
taking a valid word to the appropriate normal form. For the fully
unlabeled case use stable sorting by occupancy with label as tie breaker.
Restrict each inverse assertion to its Boolean validity tests. The
canonical codes for integer partitions are $n$ digits in base $n+1$
(use base two when $n=0$), with a positive nonincreasing prefix and
zero padding. Conjugation's inverse proof reduces coordinatewise to
$(i,j)\mapsto(j,i)$. Durfee extraction scans row indices $1$ through $n$.
\item Implement \texttt{firstHit} as least prefix length $\le2n$ at the
specified height; compare opens and closes as natural counts. The
reflection lemma states that every prefix of the flipped segment has
negated height and every later prefix changes by twice the opposite of
the old segment's height. Apply the bounded inverse-bijection count
lemma to the two failing/surplus bit masks. \texttt{firstReturn} yields
the two substring codes and lengths, hence an equivalence with pairs of
Dyck words indexed by $j$. For the partition bijection, a stack of at most
$n$ pairs (block ID, vacancies) is maintained; prefix induction verifies
that total vacancies equal opens minus closes and that finished blocks
never receive another element.
\item Tree and dissection algorithms can use bounded traversal rather than
recursion on arbitrary abstract objects. A tree with $n$ leaves and at
least two children per internal vertex has at most $n-1$ internal vertices:
the sum of all outdegrees is its number of vertices minus one and is at
least twice its number of internal vertices. Store the preorder outdegree
list, with its length, in base $2n+2$; parsing starts with one vacant
slot and replaces one slot by $d$ on reading degree $d$. The slots stay
positive before the last entry and finish at zero. A stack restores the
unique ordered tree. For plane trees with $n$ edges use $n+1$ entries
instead. Dissections can be coded as subsets of the at most $(n+2)^2$
vertex pairs on the cyclic boundary, with crossing tested by alternation
of endpoints. The base face consists of the consecutive sides incident
to its boundary walk; equivalently its vertices are exactly those outside
all open intervals cut off from the base by the selected diagonals.
This latter description is a direct bounded scan.
\item For the polynomial assertions implement \texttt{coeffMul} by the
finite convolution above and truncate after degree $N$. Prove reindexing
for pairs/triples by counting indicator arrays on $[N+1]^2$ and
$[N+1]^3$. Set \texttt{seqPos} to the sum of powers through $N$ and
prove the geometric identity since the $(N+1)$st power has all
coefficients zero through $N$. For all inductive proof invariants involving
multiple positions or sizes, fix the final input bound and verify
simultaneously all smaller indices by a Boolean fold. Signed outputs are
pairs of naturals with equality by cross-addition.
\end{itemize}
